\section{Introduction}

Wearable and portable ECG systems collect long-duration cardiac signals outside
the clinic. A static classifier trained on population data can still perform
unevenly across patients, because beat morphology, rhythm prevalence, electrode
placement, and acquisition conditions vary substantially between individuals.
The problem is most acute for minority arrhythmia classes, whose prevalence is
small relative to normal beats and whose morphology is patient-specific.

Busia \emph{et al.} proposed a 6{,}643-parameter transformer for low-power
arrhythmia classification on microcontrollers \cite{busia2025tiny}. Such a model
is attractive for embedded inference, but deployment-time personalization
introduces a distinct resource problem. Continual adaptation on the device
requires \emph{persistent} state: trainable parameters, their gradients,
optimizer moments, replay exemplars, replay labels, and buffer metadata. Under
a microcontroller-motivated memory ceiling, these components cannot be designed
independently, because every byte spent on one is a byte unavailable to the
others.

Latent replay reduces storage by preserving intermediate representations
instead of raw samples \cite{pellegrini2020latent,ravaglia2021tinyml}. In many
convolutional networks, feature-map size decreases with depth, so replaying
deeper features gives a relatively smooth memory--plasticity trade-off. The
evaluated ECG transformer behaves differently. Tokenization \emph{expands} the
signal before pooling collapses it, so the storage cost of a replay exemplar is
non-monotone in depth. We call this the \emph{replay-point cliff}: only the
post-pooling representation is smaller than the raw ECG window, and every
intermediate tap is larger.

The cliff makes an intuitive design choice look optimal: because the post-pool
vector is $\approx 11\times$ smaller than the raw beat, one can store hundreds
of post-pool exemplars and adapt only a small classifier head. That intuition is arithmetically correct about
\emph{where} replay is cheap, but it does not settle \emph{how} the budget should
be spent. Retaining a small amount of encoder plasticity through low-rank
adaptation, paired with only tens of raw replay samples, yields stronger pooled
minority-class performance and more frequent positive patient-level changes than
maximum-volume post-pool replay at the same ceiling. We do not claim it is
statistically superior: maximum-volume replay is simply \emph{not established as
uniquely optimal}, and the budget admits several competitive allocations.

This paper frames the problem as \textbf{replay--plasticity co-design} and
reports a pre-specified primary evaluation over the full eligible DS2 cohort. Relative to
an earlier exploratory version restricted to an arrhythmia-informative panel and
three seeds, every headline number here is computed over all eligible
patient-disjoint DS2 records (21 after excluding record 202) with five seeds,
with patient-bootstrap uncertainty, Holm-corrected
paired tests, and effect sizes. The contributions are:

\begin{enumerate}[leftmargin=*]
\item an architecture-level measurement of the replay-point cliff and its
consequence for where replay may be stored (Section~\ref{sec:method});
\item a complete persistent-state allocation formulation that jointly accounts
for replay representation, trainable scope, gradients, and optimizer state, and
a constrained configuration generator that turns a byte budget into a maximum
replay count (Sections~\ref{sec:method}, \ref{sec:protocol});
\item a pre-specified patient-level comparison of frozen, head-only, post-pool
adapter, and encoder-LoRA adaptation over all eligible DS2 patients with five
seeds. The metric, seed count, paired tests, effect-size definition, and multiplicity
family were fixed before the final runs. Record 202 was subsequently excluded
following a subject-identity audit, and all affected analyses were rerun without
changing the methods (Section~\ref{sec:results});
\item paired ablations separating replay location, replay volume, and trainable
depth, which support a plasticity contribution, do not support a direct replay
contribution, and find no benefit from replay volume
(Section~\ref{sec:results});
\item robustness evidence on three axes: an Adam rank-1/rank-2 budget sweep
showing that a reduced 12-record, three-seed Adam rank-1/rank-2 panel found no
statistically detectable benefit from increasing the persistent-state budget
from 8 to 64\,KiB, adaptation from
two of five trained alternative class-reweighted source checkpoints, and subject-level
cross-database transfer to INCART and SVDB (Section~\ref{sec:results});
\item a label-budget analysis quantifying how much supervision patient-specific
adaptation actually requires (Section~\ref{sec:results});
\item a memory-matched regularization baseline --- an L2 source anchor on the
$B$ factor of the low-rank delta that costs zero persistent bytes --- showing
that replay provides substantially better source retention than this tested
zero-byte $B$-factor anchor under the evaluated configuration
(Section~\ref{sec:results});
\item a leakage and reproducibility audit, and an explicit accounting of what the
software-only evidence does and does not support---including the
non-significance of the direct encoder-versus-maximum-replay comparison, the
partial coverage of the budget grid, and the small subject count on INCART
(Section~\ref{sec:limitations}).
\end{enumerate}

We deliberately avoid claiming demonstrated microcontroller deployment. The
16\,KiB ceiling is an \emph{analytical, hardware-motivated persistent-state
design constraint}: it covers what a device must retain \emph{between}
adaptation steps---trainable weights, gradients, optimizer moments, replay
exemplars, replay labels, and buffer metadata. Peak training SRAM is instead
dominated by activations (cf. TinyTL~\cite{cai2020tinytl}) and is explicitly
out of scope and unmeasured here. Runtime latency, energy, and
nonvolatile-memory endurance are likewise outside the scope of this study and
are not measured.
